% БҮЛЭГ 1 — АСУУДЛЫН ТОМЬЁОЛОЛ (УДИРТГАЛ)

\chapter{Оршил}
\label{Chapter1}

%-------------------------------------------------------------------------------
\section{Судалгааны үндэслэл}

DataReportal (2024) \cite{datareportal2024}-ийн мэдээллээр Монгол Улсад
идэвхтэй интернэт хэрэглэгчдийн эзлэх хувь нийт хүн амын 80\%-иас давж,
үүнийг даган Facebook, X, TikTok зэрэг нийгмийн сүлжээний платформ дахь
иргэдийн идэвхтэй оролцоо жилээс жилд нэмэгдэж байна. Энэхүү дижитал шилжилт
нь олон нийтийн санал бодол, хэрэглэгчийн хандлага, нийгмийн сэтгэл зүйн
төлөвийг илэрхийлэх өгөгдлийг бодит хугацаанд хуримтлуулах боломжийг нээж
буйн зэрэгцээ хувь хүний нэр төр, цахим орчны аюулгүй байдалд заналхийлсэн
доромжлол, дарамт, ташаа мэдээлэл хяналтгүй тархах хоёрдмол үр дагаврыг бий
болгож байна.

Цахим орчин дахь кибер дарамт нь хэрэглэгчдийн, ялангуяа өсвөр насныхны
сэтгэл зүйн тогтвортой байдал, нийгэмшил болон академик гүйцэтгэлд урт
хугацааны сөрөг нөлөө үзүүлдэг болохыг олон улсын судалгаануудаар тогтоосон
\cite{kowalski2014,vogels2022}. Дотоодын орчинд news.mn, gogo.mn зэрэг
мэдээллийн сайтын сэтгэгдэл болон Facebook зэрэг нийгмийн сүлжээн дэх
нийтийн хэлэлцүүлэг нь олон нийтийн байр суурийг тодорхойлох гол эх сурвалж
болж байгаа хэдий ч асуудал зөвхөн агуулгын автомат хяналтын хэрэгсэл цөөн
байгаагаар хязгаарлагдахгүй. Нийгмийн сүлжээний богино текст нь Монгол хэлний
залгамал морфологи, албан бус үгийн хэрэглээ, үсгийн алдаа, эможи, товчилсон
хэллэг, егөөдөл болон шууд илэрхийлэгдээгүй далд утгыг ихээр агуулдаг тул
стандарт бичгийн хэлд тулгуурласан загварыг шууд ашиглахад хэл шинжлэлийн
болон тооцооллын нэмэлт бэрхшээл үүсдэг.

Иймд сөрөг өнгө аястай бүх текстийг хортой агуулга гэж үзэх боломжгүй.
Хортой агуулга нь ихэвчлэн хувь хүн, бүлэг рүү чиглэсэн доромжлол, гутаалт,
дарамт, заналхийлэл зэрэг шууд халдлагын шинжтэй байдаг. Харин бүтээлч
шүүмжлэл нь тодорхой асуудал, үйлчилгээний чанар, байгууллагын үйл ажиллагаа,
эсвэл системийн доголдолд чиглэж, алдаа дутагдлыг заах, сайжруулах шаардлага
тавих, хэрэглэгчийн туршлагаа илэрхийлэх зорилготой байж болно. Мөн тухайн
хүний зан чанар бус, хийсэн ажил, шийдвэр, үйлдлийн үр дүнд чиглэсэн зөвлөгөө,
сайжруулах санал хэлбэртэй байж болно.
Ийм текст хурц үг хэллэгтэй байсан ч заавал хортой агуулга гэсэн үг биш юм.
Тиймээс автомат ангиллын систем нь зөвхөн сөрөг үгийн хэрэглээг бус, текстийн
чиглэж буй объект, илэрхийлж буй зорилго, хам сэдвийг харгалзан хортой дайралт
ба бүтээлч шүүмжлэлийг ялгах шаардлагатай.

%-------------------------------------------------------------------------------
\section{Асуудлын томьёолол}

Монгол хэл дээрх мэдрэмжийн шинжилгээний өмнөх ажлуудын дийлэнх нь текстийг
эерэг, саармаг, сөрөг гэсэн гурван ерөнхий ангиллаар ялгахад төвлөрсөн байдаг.
Энэхүү хандлага нь хэрэглэгчийн ерөнхий хандлагыг тодорхойлоход үр дүнтэй боловч
``Сөрөг'' ангиллын дотор орших хоёр өөр шинжтэй агуулгыг хангалттай
тайлбарлаж чаддаггүй. Нэг талаас хувь хүн, бүлэг рүү чиглэсэн доромжлол,
заналхийлэл, дарамт бүхий хортой сөрөг агуулга бий; нөгөө талаас үйлчилгээний
чанар, байгууллагын үйл ажиллагаа эсвэл системийн доголдолд чиглэж, сайжруулах
санал агуулсан бүтээлч шүүмжлэл оршдог.

Өөрөөр хэлбэл, хортой сөрөг агуулга ба бүтээлч шүүмжлэл хоёулаа сөрөг үг
хэллэгтэй байж болох ч ижил утгатай биш юм. Хортой сөрөг агуулга нь ихэвчлэн
хүн, бүлэг рүү чиглэсэн доромжлол, гутаалт, дарамтын шинжтэй байдаг бол
бүтээлч шүүмжлэл нь үйлчилгээ, байгууллагын шийдвэр, хийсэн ажил эсвэл
системийн доголдлыг зааж, сайжруулах шаардлага тавьдаг. Энэ ялгааг тооцохгүйгээр
бүх сөрөг текстийг нэг ангилалд оруулбал загвар бодит халдлагыг дутуу танихын
зэрэгцээ хэрэгтэй санал, асуудлын дохиог хортой агуулга гэж андуурах эрсдэлтэй.
Тиймээс Монгол кирилл текстэд хортой сөрөг агуулга ба бүтээлч шүүмжлэлийг
тусдаа ангиллаар ялгах нь энэ ажлын гол судалгааны асуудал юм.

%-------------------------------------------------------------------------------
\section{Судалгааны зорилго}

Энэхүү дипломын ажлын зорилго нь Монголын цахим орчинд нийтлэгддэг
сэтгэгдлийг эерэг, саармаг, хортой сөрөг, бүтээлч шүүмжлэл гэсэн дөрвөн
ангиллаар автоматаар ялгах эх хэлний боловсруулалтын систем боловсруулахад
оршино. Энэ хүрээнд өгөгдөл цуглуулах, цэвэрлэх, шошголох, загвар сургах,
үнэлэх болон үр дүнг хэрэглэгчид ойлгомжтой хэлбэрээр харуулах дамжлагыг нэг
системийн түвшинд нэгтгэн авч үзнэ.

Үнэлгээний хувьд зөвхөн нийт нарийвчлал өндөр байх нь хангалттай үзүүлэлт биш
юм. Дөрвөн ангиллын дундаас хортой сөрөг агуулга ба бүтээлч шүүмжлэлийн зааг
хамгийн эмзэг тул эдгээр хоёр ангиллын хооронд үүсэх андуурлыг тусад нь
шинжлэх шаардлагатай. Иймээс загваруудыг ижил өгөгдөл, ижил туршилтын нөхцөлд
харьцуулж, нийт гүйцэтгэлээс гадна ангилал тус бүрийн F1 оноо, төөрөгдлийн матриц
болон хортой сөрөг--бүтээлч шүүмжлэлийн хоорондын ялгах чадварыг хамтатган
үнэлэхийг зорив.

%-------------------------------------------------------------------------------
\section{Судалгааны зорилтууд}

Дээрх зорилгыг хэрэгжүүлэхийн тулд судалгааны ажил дараах уялдаа бүхий
зорилтуудаар тодорхойлогдоно.

\begin{enumerate}
  \item Монголын цахим орчны нийтэд нээлттэй эх сурвалжуудаас сэтгэгдлийн
        өгөгдөл цуглуулж, цэвэрлэх, давхардал арилгах, стандарт хэлбэрт
        оруулах замаар судалгааны өгөгдлийн багц бүрдүүлэх.

  \item Цуглуулсан өгөгдлийг эерэг, саармаг, хортой сөрөг, бүтээлч шүүмжлэл
        гэсэн дөрвөн ангиллаар гар аргаар шошголж, гурван шошгологчийн
        хоорондын нийцлийг Флейссийн $\kappa$ коэффициентээр хэмжин баталгаажуулах.

  \item Монгол хэлний залгамал бүтэц, кирилл бичлэгийн ялгаатай хэлбэр болон
        цахим сэтгэгдлийн албан бус найруулгыг харгалзсан урьдчилсан
        боловсруулалтын дамжлага боловсруулах.

  \item Нэг шатлалт болон хоёр шатлалт ангиллын зохион байгуулалтыг хэрэгжүүлж,
        сөрөг ангиллыг хортой сөрөг ба бүтээлч шүүмжлэл болгон нарийвчлах
        шийдлийн нөлөөг харьцуулан турших.

  \item Хэлний загварт суурилсан ангиллын хувилбарыг уламжлалт машин сургалтын
        загваруудтай ижил өгөгдөл, ижил үнэлгээний нөхцөлд харьцуулан шинжлэх.

  \item Нийт нарийвчлал, макро-F1, ангилал тус бүрийн F1 оноо болон төөрөгдлийн
        матрицад тулгуурлан загваруудын гүйцэтгэл, ялангуяа хортой сөрөг
        агуулга ба бүтээлч шүүмжлэлийн хоорондын андуурлын хэв шинжийг
        тодорхойлох.

  \item Судалгааны үр дүнг шалгаж болохуйц туршилтын прототип интерфейс
        боловсруулж, ангиллын гаралтыг хэрэглэгчид ойлгомжтой хэлбэрээр харуулах.
\end{enumerate}

%-------------------------------------------------------------------------------
\section{Таамаглал}

Судалгааны үндсэн таамаглал нь сөрөг өнгө аястай текстүүдийг нэг ангилалд
нэгтгэх үед хортой сөрөг агуулга ба бүтээлч шүүмжлэлийн хоорондын ялгаа
бүдгэрнэ гэсэн санаанд тулгуурлана. Иймээс энэ ажилд зөвхөн эерэг, саармаг,
сөрөг гэсэн ерөнхий ялгалтаар хязгаарлахгүйгээр сөрөг агуулгын дотоод ялгааг
тусад нь авч үзэж, өгөгдлийн чанар болон ангиллын зохион байгуулалтын нөлөөг
дараах дэд таамаглалаар шалгана.

\begin{itemize}
  \item H$_1$: Хортой сөрөг агуулга ба бүтээлч шүүмжлэлийг тусдаа ангиллаар
    авч үзэх нь ерөнхий сөрөг ангилалд далдлагдах андуурлын хэв шинжийг илүү
    тодорхой харуулна.
  \item H$_2$: Өгөгдөл цэвэрлэх, давхардал арилгах, маргаантай жишээг дахин
    нягтлах зэрэг өгөгдлийн чанаржуулалтын алхам нь нийт нарийвчлалаас гадна
    макро-F1 болон ангилал тус бүрийн F1 оноонд мэдэгдэхүйц нөлөө үзүүлнэ.
  \item H$_3$: Хоёр шатлалт зохион байгуулалт нь хортой сөрөг ба бүтээлч
    шүүмжлэлийг тусгайлан ялгах тайлбарлагдахуйц бүтэц санал болгох боловч
    шат дамжсан алдааны улмаас нийт системийн
    гүйцэтгэл буурах эрсдэлтэй.
\end{itemize}

%-------------------------------------------------------------------------------
\section{Судалгааны хамрах хүрээ ба хязгаарлалт}

Судалгааны хамрах хүрээ нь Монгол кирилл бичгийн цахим текстээр хязгаарлагдана.
Латин галиг, хэл холилдсон текст болон нийгмийн сүлжээнд өргөн тархсан
товчилсон бичлэгийн хэлбэрүүд нь MN-BERT токенчлолын тогтвортой байдал болон
шошгололтын нэг утгатай байдлыг алдагдуулах эрсдэлтэй тул энэ ажлын өгөгдлийн
хүрээнд оруулаагүй. Өгөгдөл нь news.mn, gogo.mn болон Facebook зэрэг нийгмийн
сүлжээний нийтэд нээлттэй эх сурвалжид тулгуурлах бөгөөд 2025--2026 оны
хугацаанд цуглуулсан текстэн өгөгдлөөр хязгаарлагдана.

Мөн боловсруулсан систем нь үйлдвэрлэлийн орчинд шууд байршуулсан эцсийн
бүтээгдэхүүн бус, судалгааны үр дүнг шалгах зориулалттай прототип шийдэл юм.
Иймээс бодит хэрэглээнд нэвтрүүлэх тохиолдолд өгөгдлийн өргөтгөл, хүний
хяналттай шийдвэр гаргалт, тайлбарлагдахуйц байдлын нэмэлт механизм шаардлагатай
гэж үзэв.

%-------------------------------------------------------------------------------
\section{Ажлын шинэлэг тал ба ач холбогдол}

Энэхүү ажлын шинэлэг тал нь Монгол цахим орчны кирилл сэтгэгдэлд хувь хүн,
бүлэг рүү чиглэсэн хортой сөрөг агуулгыг үйлчилгээ, байгууллагын үйл ажиллагаа,
хийсэн ажил, шийдвэрт чиглэж сайжруулах санал агуулсан бүтээлч шүүмжлэлээс
тусдаа ангиллаар ялгасанд оршино.

Судалгааны гол хувь нэмэр нь дараах гурван чиглэлээр тодорхойлогдоно.
Нэгдүгээрт, Монгол кирилл цахим сэтгэгдэлд хортой сөрөг агуулга ба бүтээлч
шүүмжлэлийг тусдаа ангиллаар авч үзэх ангиллын хүрээг тодорхойлсон.
Хоёрдугаарт, news.mn, gogo.mn болон Facebook зэрэг нийтэд нээлттэй эх сурвалжаас
цуглуулсан өгөгдлийг дөрвөн ангиллаар шошголж, Флейссийн $\kappa = 0.72$
үзүүлэлтээр шошгологчдын нийцлийг баталгаажуулсан өгөгдлийн багц бүрдүүлсэн.
Гуравдугаарт, MN-BERT-д суурилсан нэг шатлалт болон хоёр шатлалт ангиллын
зохион байгуулалтыг уламжлалт машин сургалтын загваруудтай харьцуулан туршиж,
хортой сөрөг ба бүтээлч шүүмжлэлийн зааг дээр үүсэх алдааны хэв шинжийг
шинжилсэн.

Энэхүү ажил нь Монгол кирилл текстийн автомат ангиллын цаашдын судалгаа болон
хэрэглээний суурь болно.
